\sloppy
\emergencystretch=5em
\hbadness=10000
\vbadness=10000
\tolerance=2000
\providecolor{codebg}{RGB}{248,248,248}
\providecolor{codeframe}{RGB}{200,200,200}
\providecolor{keyword}{RGB}{0,0,180}
\providecolor{string}{RGB}{163,21,21}
\providecolor{comment}{RGB}{0,128,0}
\lstset{
	basicstyle=\ttfamily\scriptsize,
	breaklines=true,
	breakatwhitespace=false,
	columns=fullflexible,
	keepspaces=true,
	frame=single,
	framesep=2pt,
	showstringspaces=false,
	backgroundcolor=\color{codebg},
	rulecolor=\color{codeframe},
	keywordstyle=\color{keyword}\bfseries,
	stringstyle=\color{string},
	commentstyle=\color{comment}\itshape,
	numberstyle=\tiny\color{gray},
	numbers=left,
	numbersep=6pt,
	xleftmargin=14pt,
	framexleftmargin=14pt,
	literate=
	{ą}{{\v{a}}}1 {ę}{{\k{e}}}1 {ż}{{\.z}}1 {ś}{{\'s}}1
	{ł}{{\l}}1 {ć}{{\'c}}1 {ń}{{\'n}}1 {ó}{{\'o}}1 {ź}{{\'z}}1
	{→}{{$\rightarrow$}}1 {←}{{$\leftarrow$}}1 {≠}{{$\neq$}}1
}

\section{نمای کلی و هدف ماژول}

فایل \lr{\texttt{qkd/proofs/optimization.py}} مجموعه‌ای از ابزارهای بهینه‌سازی پارامتر برای پروتکل‌های مختلف توزیع کلید کوانتومی (\lr{QKD}) را پیاده‌سازی می‌کند. این ماژول بر اساس دو چارچوب مرجع توسعه یافته است: چارچوب \lr{Ma et al. 2005} برای پروتکل‌های \lr{Decoy-State BB84} و چارچوب \lr{Chapman et al. 2018} برای پروتکل‌های \lr{state-dependent} مبتنی بر \lr{Amplitude Damping Channel (ADC)}. هدف اصلی این ماژول، پیدا کردن مقادیر بهینۀ پارامترهای قابل کنترل آلیس (مانند شدت سیگنال $\mu$، شدت \lr{decoy} $\nu$، احتمال انتخاب سیگنال $P_s$، احتمال انتخاب \lr{decoy} $P_d$، و زاویۀ چرخش کدگذاری $\theta$) به گونه‌ای است که نرخ کلید امن نهایی به حداکثر برسد. این بهینه‌سازی در عمل برای استقرار سیستم‌های \lr{QKD} حیاتی است، زیرا انتخاب پارامترهای غیربهینه می‌تواند نرخ کلید را تا چندین براکم کند یا حتی آن را به صفر برساند.

اهمیت این ماژول از این واقعیت ناشی می‌شود که نرخ کلید امن $R$ در پروتکل‌های \lr{QKD} یک تابع غیرخطی و معمولاً غیرمحدب (\lr{non-convex}) از پارامترهای سیستم است. به عنوان مثال، در پروتکل \lr{Decoy-State BB84}، نرخ کلید به صورت $R = q\{Q_1[1 - H_2(e_1)] - Q_\mu f(E_\mu) H_2(E_\mu)\}$ محاسبه می‌شود که در آن $Q_1 = Y_1 \mu e^{-\mu}$ و $Q_\mu = \sum_n Y_n P(n|\mu)$ هر دو به صورت غیرخطی به $\mu$ وابسته‌اند. افزایش $\mu$ باعث افزایش $Q_1$ (و در نتیجه افزایش ترم تقویت حریم خصوصی) می‌شود، اما همزمان باعث افزایش $Q_\mu$ (و در نتیجه افزایش نشت تصحیح خطا) و افزایش آسیب‌پذیری به حمله \lr{PNS} نیز می‌شود. بنابراین یک نقطۀ بهینۀ $\mu^*$ وجود دارد که در آن توازن میان این دو ترم بهینه است. پیدا کردن این نقطه به صورت تحلیلی به دلیل پیچیدگی فرمول‌ها ممکن نیست و باید از روش‌های عددی استفاده شود.

این ماژول ده تابع بهینه‌سازی مستقل ارائه می‌کند که هر کدام برای یک سناریوی خاص طراحی شده‌اند. این توابع عبارتند از: \lr{\texttt{optimize\_decoy\_intensities}} (بهینه‌سازی اکتشافی شدت سیگنال و \lr{decoy} بر اساس مدل نرخ ساده‌شده)، \lr{\texttt{optimize\_rotation\_angle}} (بهینه‌سازی زاویۀ چرخش کدگذاری برای یک کانال عمومی)، \lr{\texttt{calculate\_optimal\_mu\_ma2005}} (محاسبۀ تحلیلی-عددی $\mu$ بهینه بر اساس معادلۀ 12 مقالۀ \lr{Ma 2005})، \lr{\texttt{optimize\_nu\_ma2005}} (بهینه‌سازی $\nu$ برای $N$ متناهی بر اساس بخش 4.3 مقاله)، \lr{\texttt{optimize\_pulse\_distribution\_ma2005}} (بهینه‌سازی احتمالات توزیع پالس با روش \lr{SLSQP})، \lr{\texttt{optimize\_pulse\_probabilities}} (پوشش‌دهنده برای تابع قبلی)، \lr{\texttt{optimize\_simultaneous\_ma2005}} (بهینه‌سازی همزمان $P_s$، $P_d$ و $\nu$ با استفاده از کلاس \lr{\texttt{Ma2005GlobalOptimizer}} و گرادیان عددی)، \lr{\texttt{optimize\_global\_ma2005}} (پوشش‌دهنده برای تابع قبلی)، \lr{\texttt{optimize\_lim2014\_finite\_key}} (بهینه‌سازی پنج پارامتر چارچوب \lr{Lim 2014} برای کلید متناهی با روش \lr{Nelder-Mead})، و \lr{\texttt{optimize\_rotation\_angle\_chapman}} (بهینه‌سازی زاویۀ چرخش برای \lr{Coherent Scheme} در \lr{Chapman 2018} با محاسبۀ فاصلۀ \lr{trace distance} میان حالت‌های خروجی کانال).

طراحی این ماژول بر اساس اصل جداسازی دغدغه‌ها (\lr{Separation of Concerns}) استوار است؛ به این ترتیب که توابع بهینه‌سازی صرفاً مسئول پیدا کردن پارامترهای بهینه هستند و محاسبۀ نرخ کلید (یا هر کمیت هدف دیگر) به یک فراخوان (\lr{callback}) تابعی سپرده می‌شود که توسط کاربر ارائه می‌گردد. این طراحی چندین مزیت دارد: 

نخست، ماژول بهینه‌سازی از جزئیات پیاده‌سازی اثبات امنیتی مستقل است و می‌تواند با هر چارچوب اثبات (مانند \lr{Ma 2005}، \lr{Lim 2014} یا \lr{Tight}) کار کند؛ دوم، کاربر می‌تواند تابع هدف خود را به سادگی جایگزین کند (مثلاً به جای نرخ کلید، \lr{QBER} را کمینه کند یا طول کلید متناهی را بیشینه سازد) بدون آنکه نیازی به تغییر کد بهینه‌سازی باشد؛ سوم، فرآیند تست‌کردن و ارزیابی آسان‌تر می‌شود، زیرا می‌توان از \lr{callback}‌های مصنوعی با رفتار مشخص و شناخته‌شده برای اعتبارسنجی توابع بهینه‌سازی استفاده کرد. این الگو که به تزریق وابستگی (\lr{«dependency injection»}) نیز معروف است، یک اصل کلیدی در مهندسی نرم‌افزار به شمار می‌رود که قابلیت نگهداری و توسعهٔ کد را به طور قابل‌توجهی افزایش می‌دهد.


در صورت شکست بهینه‌سازی (به هر دلیل: عدم همگرایی، ناحیه غیرممکن، یا خطای عددی)، توابع این ماژول به جای \lr{raise} کردن استثنا، مقادیر پیش‌فرض امن و محافظه‌کارانه را برمی‌گردانند. این رفتار \lr{fail-safe} تضمین می‌کند که خط لوله شبیه‌سازی متوقف نشود و کاربر بتواند با مقادیر پیش‌فرض ادامه دهد. به عنوان مثال، اگر بهینه‌سازی $\mu$ شکست بخورد، مقدار پیش‌فرض 0.48 (برای $f \le 1.05$) یا 0.54 (برای $f > 1.05$) برگردانده می‌شود که مقادیر تجربی خوبی برای پروتکل‌های \lr{Decoy-State BB84} هستند. این مقادیر پیش‌فرض بر اساس بخش 3.1 مقالۀ \lr{Ma 2005} انتخاب شده‌اند.

\section{مبانی تئوری و فیزیکی}

\subsection{نرخ کلید \lr{GLLP} و وابستگی غیرخطی به $\mu$}

فرمول استاندارد نرخ کلید در چارچوب \lr{Gottesman-Lo-Lütkenhaus-Preskill (GLLP)} به صورت زیر است:
\begin{equation}
	R = q\left\{Q_1\left[1 - H_2(e_1)\right] - Q_\mu\,f(E_\mu)\,H_2(E_\mu)\right\},
\end{equation}
که در آن $q$ فاکتور پایه‌گذاری (برای \lr{BB84} برابر \lr{0.5})، $Q_1 = Y_1 \mu e^{-\mu}$ \lr{Gain} تک‌فوتونی، $Q_\mu = \sum_n Y_n P(n|\mu)$ \lr{Gain} کلی سیگنال، $e_1$ خطای تک‌فوتونی، $E_\mu$ \lr{QBER} سیگنال، $f(E_\mu)$ بازدهی تصحیح خطا و $H_2(\cdot)$ آنتروپی باینری است. ترم اول $Q_1[1 - H_2(e_1)]$ نشان‌دهندۀ آنتروفی قابل استخراج از تک‌فوتون‌های امن است و ترم دوم $Q_\mu f(E_\mu) H_2(E_\mu)$ نشان‌دهندۀ نشت تصحیح خطا است. 

برای پیدا کردن $\mu$ بهینه، باید $\partial R / \partial \mu = 0$ را حل کنیم. با فرض ساده‌کنندۀ $e_1 \approx E_\mu$ و $Y_1 \approx \eta$ (در حد نویز کم)، این مشتق‌گیری به یک معادلۀ تحلیلی-عددی منجر می‌شود که در معادلۀ (12) مقالۀ \lr{Ma 2005} ارائه شده است. این معادله به صورت زیر فرمول‌بندی می‌شود:
\begin{equation}
	(1 - \mu)\,e^{-\mu} = C, \qquad C = \frac{f_{\text{EC}}\,H_2(e_{\text{det}})}{1 - H_2(e_{\text{det}})},
\end{equation}
که در آن $e_{\text{det}}$ خطای ذاتی آشکارساز است. تابع $(1-\mu)e^{-\mu}$ در $\mu = 0$ برابر 1 است و به طور یکنواخت تا $\mu \to \infty$ کاهش می‌یابد. بنابراین برای هر $C \in (0, 1)$ یک جواب یکتا در بازۀ $(0, 1)$ وجود دارد. اگر $C \ge 1$ (یعنی خطای آشکارساز خیلی بزرگ باشد)، جوابی وجود ندارد و نرخ کلید منفی می‌شود؛ در این حالت، ماژول به مقدار پیش‌فرض $\mu = 0.48$ برمی‌گردد. این محاسبۀ تحلیلی-عددی در تابع \lr{\texttt{calculate\_optimal\_mu\_ma2005}} با استفاده از روش \lr{Brent's method} (\lr{\texttt{brentq}} از \lr{SciPy}) پیاده‌سازی شده است.

\subsection{بهینه‌سازی $\nu$ برای کلید متناهی}

در حالت کلید متناهی ($N < \infty$)، نرخ کلید به اندازۀ نمونۀ آماری \lr{decoy} حساس می‌شود. اگر $\nu$ خیلی کوچک باشد، تعداد آشکارسازی‌های \lr{decoy} $N_d = N \cdot P_d \cdot Q_\nu$ کم می‌شود و کران‌های آماری برای $Y_1$ و $e_1$ سخت‌گیرانه‌تر (یعنی بدتر) می‌شوند. اگر $\nu$ خیلی بزرگ باشد، شباهت آن به $\mu$ زیاد می‌شود و قابلیت تفکیک آماری بین سیگنال و \lr{decoy} کاهش می‌یابد. بنابراین یک $\nu^*$ بهینه وجود دارد. بخش 4.3 مقالۀ \lr{Ma 2005} این بهینه‌سازی را به صورت یک مینیمم‌سازی اسکالر روی $\nu$ با کران‌های $[\nu_{\min}, \mu \cdot 0.9]$ فرمول‌بندی می‌کند. در کد، این بهینه‌سازی در تابع \lr{\texttt{optimize\_nu\_ma2005}} با استفاده از روش \lr{Brent's bounded method} (\lr{\texttt{minimize\_scalar}} با \lr{\texttt{method='bounded'}}) پیاده‌سازی شده است. این روش برای توابع یک‌بعدی هموار بهینه است و همگرایی تضمین‌شده‌ای در بازۀ مشخص‌شده دارد.

\subsection{بهینه‌سازی توامان $P_s$، $P_d$ و $\nu$}

بهینه‌سازی توامان سه پارامتر $P_s$، $P_d$ و $\nu$ یک مسألۀ بهینه‌سازی چندبعدی است که نیاز به روش‌های \lr{gradient-based} دارد. در کد، این مسأله با استفاده از روش \lr{SLSQP (Sequential Least SQuares Programming)} از \lr{SciPy} حل می‌شود. این روش برای مسائل بهینه‌سازی مقید با محدودیت‌های نامساوی مناسب است و از گرادیان تابع هدف برای هدایت جستجو استفاده می‌کند. کلاس \lr{\texttt{Ma2005GlobalOptimizer}} برای این منظور طراحی شده و یک متد \lr{\texttt{gradient}} برای محاسبۀ گرادیان به روش تفاضل محدود (\lr{finite difference}) ارائه می‌کند. 

گرادیان به روش تفاضل محدود مرکزی به صورت زیر محاسبه می‌شود:
\begin{equation}
	\frac{\partial f}{\partial x_i} \approx \frac{f(x + \epsilon\,e_i) - f(x - \epsilon\,e_i)}{2\,\epsilon},
\end{equation}
که در آن $e_i$ بردار واحد در جهت $i$-ام و $\epsilon$ گام تفاضل محدود است (در کد $\epsilon = 10^{-8}$). این روش برای $P_s$ و $P_d$ به صورت تفاضل پیشرو (\lr{forward difference}) و برای $\nu$ به صورت تفاضل مرکزی پیاده‌سازی شده است. این انتخاب به این دلیل است که $P_s$ و $P_d$ در محدودیت‌های سخت ($P_s + P_d \le 1$) شرکت می‌کنند و تفاضل پیشرو در نزدیک مرز محدودیت پایدارتر است، در حالی که $\nu$ یک متغیر آزاد است و تفاضل مرکزی دقت بالاتری دارد.

\subsection{کانال \lr{Amplitude Damping} و \lr{Chapman 2018}}

در پروتکل‌های \lr{state-dependent QKD} که توسط \lr{Chapman et al. 2018} پیشنهاد شده‌اند، آلیس اطلاعات را در حالت‌های کوانتومی $\rho_0$ و $\rho_1$ کدگذاری می‌کند و این حالت‌ها از یک کانال \lr{Amplitude Damping Channel (ADC)} با پارامتر $\gamma$ عبور می‌کنند. \lr{ADC} کانالی است که افت انرژی (مانند واهلش آرام اتم‌ها) را مدل می‌کند و عملگرهای \lr{Kraus} آن به صورت زیر تعریف می‌شوند:
\begin{align}
	E_0 &= \begin{pmatrix} 1 & 0 \\ 0 & \sqrt{1 - \gamma} \end{pmatrix}, \\
	E_1 &= \begin{pmatrix} 0 & \sqrt{\gamma} \\ 0 & 0 \end{pmatrix}.
\end{align}
حالت خروجی کانال از رابطۀ $\rho_{\text{out}} = E_0 \rho_{\text{in}} E_0^\dagger + E_1 \rho_{\text{in}} E_1^\dagger$ محاسبه می‌شود. در \lr{Coherent Scheme}، آلیس حالت‌های ورودی $|0\rangle\langle 0|$ و $|1\rangle\langle 1|$ را با یک چرخش \lr{Y} به زاویۀ $\theta$ تبدیل می‌کند:
\begin{equation}
	\rho_i^{\text{in}}(\theta) = R_y(\theta)\,|i\rangle\langle i|\,R_y(\theta)^\dagger, \qquad R_y(\theta) = \begin{pmatrix} \cos(\theta/2) & -\sin(\theta/2) \\ \sin(\theta/2) & \cos(\theta/2) \end{pmatrix}.
\end{equation}
احتمال موفقیت تشخیص حالت توسط باب متناسب با فاصلۀ \lr{trace distance} میان $\rho_0^{\text{out}}$ و $\rho_1^{\text{out}}$ است:
\begin{equation}
	P_{\text{succ}} = \frac{1}{2}\left(1 + D(\rho_0^{\text{out}}, \rho_1^{\text{out}})\right), \qquad D(\rho_0, \rho_1) = \frac{1}{2}\,\text{Tr}\left|\rho_0 - \rho_1\right|,
\end{equation}
که در آن $|A| = \sqrt{A^\dagger A}$ و \lr{Tr} عملگر \lr{trace} است. برای ماتریس‌های \lr{Hermitian}، $D$ را می‌توان به صورت $D = \frac{1}{2}\sum_i |\lambda_i|$ محاسبه کرد که در آن $\lambda_i$ مقادیر ویژۀ $\rho_0 - \rho_1$ هستند. بهینه‌سازی زاویۀ $\theta$ به گونه‌ای که $D$ بیشینه شود، در تابع \lr{\texttt{optimize\_rotation\_angle\_chapman}} پیاده‌سازی شده است.

\subsection{بهینه‌سازی کلید متناهی \lr{Lim 2014}}

چارچوب \lr{Lim et al. 2014} پنج پارامتر آزاد دارد: $q_x$ (احتمال پایۀ $X$)، $P_{\mu_1}$ (احتمال سیگنال)، $P_{\mu_2}$ (احتمال \lr{decoy} ضعیف)، $\mu_1$ (شدت سیگنال) و $\mu_2$ (شدت \lr{decoy}). بهینه‌سازی نرخ کلید $R = \ell / N$ روی این پنج پارامتر یک مسألۀ پیچیده است زیرا چشم‌انداز (\lr{landscape}) تابع هدف معمولاً ناهموار (\lr{non-smooth}) است و دارای چندین محلی است. در کد، این مسأله با روش \lr{Nelder-Mead} (یک روش بدون گرادیان) حل می‌شود که برای توابع ناهموار مقاوم‌تر از روش‌های مبتنی بر گرادیان است. این روش از یک \lr{simplex} (چندوجهی محدب) برای کاوش فضای پارامتر استفاده می‌کند و با منبسط‌کردن، منقبض‌کردن و انعکاس دادن \lr{simplex}، به سمت نقطۀ بهینه حرکت می‌کند. حداکثر تعداد تکرارها در کد 200 تنظیم شده و تلورانس $10^{-3}$ برای توقف استفاده می‌شود.

\subsection{محدودیت‌های فیزیکی و منطق جریمه}

در تمام توابع بهینه‌سازی این ماژول، محدودیت‌های فیزیکی به صورت جریمه‌های (\lr{penalty}) بزرگ (مانند $10^9$) پیاده‌سازی شده‌اند. این محدودیت‌ها عبارتند از: $\mu > \nu$ (شدت سیگنال باید بزرگتر از \lr{decoy} باشد)، $\nu > 0$ (شدت باید مثبت باشد)، $\mu \le 1$ (برای جلوگیری از پالس‌های بسیار چندفوتونی)، $P_s + P_d < 1$ (احتمالات باید نرمالیزه شوند و فضایی برای $P_v$ باقی بماند)، و $\nu < 0.9 \mu$ (برای حفظ قابلیت تفکیک آماری). استفاده از جریمه به جای محدودیت‌های سخت به روش \lr{Nelder-Mead} اجازه می‌دهد در ناحیه‌های غیرممکن هم کاوش کند (با جریمه بالا) و در عین حال به روش \lr{SLSQP} (که از محدودیت‌های صریح پشتیبانی می‌کند) انعطاف‌پذیری بدهد. این طراحی ترکیبی به ماژول اجازه می‌دهد با انواع روش‌های بهینه‌سازی کار کند بدون آنکه نیاز به بازنویسی تابع هدف باشد.

\section{تحلیل معماری و ساختار داده‌ها}

\subsection{ساختار کلی ماژول}

ماژول \lr{\texttt{optimization.py}} بر اساس یک طراحی تابعی (\lr{functional design}) استوار است که در آن نه توابع مستقل بهینه‌سازی ارائه می‌شوند و یک کلاس کمکی \lr{\texttt{Ma2005GlobalOptimizer}} برای بهینه‌سازی توامان استفاده می‌شود. این طراحی بر اساس اصل \lr{``functions first, classes when needed''} است که در آن توابع ساده برای کاربردهای یک‌بار مصرف و کلاس‌ها برای حالت‌هایی که نیاز به نگهداری \lr{state} (مانند \lr{callback} و گام تفاضل محدود) هستند، استفاده می‌شوند. تمام توابع در لیست \lr{\texttt{\_\_all\_\_}} صریحاً اعلام شده‌اند تا به کاربر یک رابط واضح از API عمومی ماژول ارائه دهند.

\subsection{کلاس \lr{Ma2005GlobalOptimizer}}

کلاس \lr{\texttt{Ma2005GlobalOptimizer}} برای بهینه‌سازی توامان $P_s$، $P_d$ و $\nu$ طراحی شده است. در ادامه، ساختار این کلاس را به صورت خلاصه نمایش می‌دهیم:

\begin{LTR}
	\begin{lstlisting}[language=Python]
		class Ma2005GlobalOptimizer:
		"""
		Implements the global optimization for Vacuum+Weak decoy protocol
		using analytical derivatives derived from Eqs. (34) and (37) w.r.t nu
		and Ps, Pd.
		"""
		def __init__(self, mu_target, n_total, cost_callback):
		self.mu = mu_target
		self.N = n_total
		self.callback = cost_callback
		self.epsilon = 1e-8  # small finite diff step
	\end{lstlisting}
\end{LTR}

سازندۀ این کلاس چهار پارامتر دریافت می‌کند: \lr{\texttt{mu\_target}} (شدت هدف سیگنال $\mu$)، \lr{\texttt{n\_total}} (تعداد کل پالس‌های $N$)، \lr{\texttt{cost\_callback}} (تابع هدف که نرخ منفی کلید را محاسبه می‌کند) و \lr{\texttt{epsilon}} (گام تفاضل محدود با مقدار پیش‌فرض $10^{-8}$). این کلاس به صورت یک \lr{stateful wrapper} عمل می‌کند که \lr{callback} و پارامترهای سیستم را نگه می‌دارد و متدهای \lr{\texttt{objective}} و \lr{\texttt{gradient}} را برای استفاده توسط \lr{solver} \lr{SciPy} ارائه می‌دهد. این طراحی از الگوی \lr{``Strategy''} پیروی می‌کند که در آن الگوریتم بهینه‌سازی (\lr{SLSQP}) از تابع هدف (که توسط \lr{callback} تعریف می‌شود) مستقل است.

\subsection{متد \lr{\texttt{objective}} و جریمه‌ها}

متد \lr{\texttt{objective}} تابع هدف را برای \lr{solver} فراهم می‌کند:

\begin{LTR}
	\begin{lstlisting}[language=Python]
		def objective(self, x):
		"""Objective function for minimization (-Rate)."""
		p_s, p_d, nu = x
		if p_s <= 0.01 or p_d <= 0.01 or (p_s + p_d) >= 1.0:
		return 1e9
		if nu <= 0.001 or nu >= self.mu * 0.99:
		return 1e9
		# Call the provided cost function (which computes -Rate)
		return self.callback(p_s, p_d, nu)
	\end{lstlisting}
\end{LTR}

در این متد، ابتدا محدودیت‌های فیزیکی بررسی می‌شوند: $P_s > 0.01$، $P_d > 0.01$، $P_s + P_d < 1$، $\nu > 0.001$ و $\nu < 0.99 \mu$. اگر هر کدام از این محدودیت‌ها نقض شود، یک جریمۀ بزرگ ($10^9$) برگردانده می‌شود. این جریمه به \lr{solver} می‌گوید که این ناحیه از فضای پارامتر نامطلوب است و باید از آن دوری کند. انتخاب کران‌های 0.01 و 0.001 به این دلیل است که مقادیر خیلی کوچک از نظر آماری معنایی ندارند (تعداد پالس‌های \lr{decoy} خیلی کم می‌شود) و از نظر عددی نیز ناپایدار هستند. اگر محدودیت‌ها نقض نشوند، \lr{callback} فراخوانی می‌شود و مقدار تابع هدف (نرخ منفی کلید) برگردانده می‌شود. این طراحی به کاربر اجازه می‌دهد تابع هدف خود را بدون نگرانی در مورد محدودیت‌های فیزیکی پیاده‌سازی کند.

\subsection{متد \lr{\texttt{gradient}} با تفاضل محدود}

متد \lr{\texttt{gradient}} گرادیان تابع هدف را به روش تفاضل محدود محاسبه می‌کند:

\begin{LTR}
	\begin{lstlisting}[language=Python]
		def gradient(self, x):
		"""Calculates the gradient [d(-R)/dPs, d(-R)/dPd, d(-R)/dnu]."""
		p_s, p_d, nu = x
		
		# 1. Calculate current rate
		val_center = self.objective(x)
		if val_center >= 1e8:
		return np.zeros(3)  # Infeasible region
		
		# 2. Finite Difference for Probabilities (Ps, Pd)
		grad = np.zeros(3)
		
		# d/dPs (forward difference)
		x_ps = x.copy(); x_ps[0] += self.epsilon
		val_ps = self.objective(x_ps)
		grad[0] = (val_ps - val_center) / self.epsilon
		
		# d/dPd (forward difference)
		x_pd = x.copy(); x_pd[1] += self.epsilon
		val_pd = self.objective(x_pd)
		grad[1] = (val_pd - val_center) / self.epsilon
		
		# 3. Centered difference for nu
		x_nu_p = x.copy(); x_nu_p[2] += self.epsilon
		x_nu_m = x.copy(); x_nu_m[2] -= self.epsilon
		grad[2] = (self.objective(x_nu_p) - self.objective(x_nu_m)) / (2 * self.epsilon)
		
		return grad
	\end{lstlisting}
\end{LTR}

این متد گرادیان را به صورت یک بردار سه‌بعدی $[\partial(-R)/\partial P_s, \partial(-R)/\partial P_d, \partial(-R)/\partial \nu]$ محاسبه می‌کند. ابتدا مقدار تابع هدف در نقطۀ فعلی محاسبه می‌شود. اگر این مقدار در ناحیۀ غیرممکن باشد (بزرگتر از $10^8$)، گرادیان صفر برگردانده می‌شود تا \lr{solver} در همان نقطه بماند. سپس برای $P_s$ و $P_d$ از تفاضل پیشرو استفاده می‌شود: $\partial f / \partial x_i \approx [f(x + \epsilon e_i) - f(x)] / \epsilon$. این روش به یک ارزیابی تابع اضافی برای هر پارامتر نیاز دارد (در مجموع دو ارزیابی برای $P_s$ و $P_d$). برای $\nu$ از تفاضل مرکزی استفاده می‌شود: $\partial f / \partial \nu \approx [f(\nu + \epsilon) - f(\nu - \epsilon)] / (2\epsilon)$. این روش به دو ارزیابی اضافی نیاز دارد اما دقت مرتبۀ دوم (\lr{second-order accuracy}) دارد (خطای $O(\epsilon^2)$ در مقابل $O(\epsilon)$ برای تفاضل پیشرو). انتخاب ترکیبی تفاضل پیشرو برای $P_s$ و $P_d$ و تفاضل مرکزی برای $\nu$ به این دلیل است که $P_s$ و $P_d$ در محدودیت‌های سخت شرکت می‌کنند و تفاضل پیشرو در نزدیک مرز پایدارتر است، در حالی که $\nu$ یک متغیر آزاد است و تفاضل مرکزی دقت بالاتری ارائه می‌دهد.

\subsection{الگوی \lr{callback} و تزریق وابستگی}

یکی از ویژگی‌های کلیدی این ماژول، استفاده گسترده از الگوی \lr{callback} است. توابع بهینه‌سازی به جای آنکه خود نرخ کلید را محاسبه کنند، یک تابع \lr{callback} از کاربر دریافت می‌کنند که این محاسبه را انجام می‌دهد. به عنوان مثال، تابع \lr{\texttt{optimize\_simultaneous\_ma2005}} یک \lr{callback} با امضای \lr{\texttt{Callable[[float, float, float], float]}} دریافت می‌کند که به این معناست تابعی است که سه عدد اعشاری ($P_s$، $P_d$، $\nu$) دریافت می‌کند و یک عدد اعشاری (نرخ منفی کلید) برمی‌گرداند. این طراحی چندین مزیت دارد:

\begin{itemize}
	\item جداسازی لایه‌ها: لایۀ بهینه‌سازی از لایۀ اثبات امنیتی مستقل است.
	\item انعطاف‌پذیری: کاربر می‌تواند تابع هدف خود را به‌سادگی جایگزین کند (مثلاً به‌جای نرخ کلید، \lr{QBER} را کمینه کند).
	\item تست‌پذیری: می‌توان \lr{callback}‌های مصنوعی با رفتار شناخته‌شده برای اعتبارسنجی استفاده کرد.
	\item قابلیت توسعه: اضافه‌کردن چارچوب‌های اثبات جدید نیازی به تغییر کد بهینه‌سازی ندارد.
\end{itemize}

این الگو که به \lr{``dependency injection''} نیز شناخته می‌شود، یک اصل کلیدی مهندسی نرم‌افزار است که قابلیت نگهداری و توسعه کد را به طور قابل توجهی افزایش می‌دهد.

\subsection{Type Hint ها و ایمنی نوع}

ماژول به طور گسترده از \lr{Type Hint} های \lr{Python} استفاده می‌کند تا ایمنی نوع و قابلیت خواندن کد را افزایش دهد. به عنوان مثال:

\begin{LTR}
	\begin{lstlisting}[language=Python]
		def optimize_simultaneous_ma2005(
		mu_target: float,
		total_pulses_N: int,
		cost_function_callback: Callable[[float, float, float], float]
		) -> Tuple[float, float, float, float]:
	\end{lstlisting}
\end{LTR}

در این امضا، \lr{\texttt{mu\_target: float}} مشخص می‌کند که \lr{\texttt{mu\_target}} باید یک عدد اعشاری باشد، \lr{\texttt{total\_pulses\_N: int}} مشخص می‌کند که \lr{\texttt{total\_pulses\_N}} باید یک عدد صحیح باشد، و \lr{\texttt{cost\_function\_callback: Callable[[float, float, float], float]}} یک \lr{type hint} پیچیده‌تر است که مشخص می‌کند \lr{callback} یک تابع است که سه عدد اعشاری دریافت می‌کند و یک عدد اعشاری برمی‌گرداند. خروجی تابع نیز به صورت \lr{\texttt{Tuple[float, float, float, float]}} مشخص شده که نشان می‌دهد یک تاپل چهار عنصری از اعداد اعشاری برگردانده می‌شود ($P_s^*$، $P_d^*$، $P_v^*$، $\nu^*$). این \lr{Type Hint} ها به ویرایشگرهای کد اجازه می‌دهند خطاهای نوع را در زمان توسعه شناسایی کنند و به کاربر در فهم رابط تابع کمک می‌کنند.

\section{پیاده‌سازی ریاضی و الگوریتمی}

\subsection{محاسبۀ $\mu$ بهینه بر اساس معادلۀ 12 مقالۀ \lr{Ma 2005}}

تابع \lr{\texttt{calculate\_optimal\_mu\_ma2005}} محاسبۀ تحلیلی-عددی $\mu$ بهینه را پیاده‌سازی می‌کند. در ادامه، نسخۀ کامل این تابع را با توضیح خط به خط ارائه می‌دهیم:

\begin{LTR}
	\begin{lstlisting}[language=Python]
		def calculate_optimal_mu_ma2005(
		distance_km: float,
		fiber_loss_db_km: float,
		det_eff: float,
		dark_rate: float,
		error_correction_f: float = 1.22,
		e_detector: float = 0.033
		) -> float:
		"""
		Calculates optimal signal intensity (mu) maximizing the GLLP rate.
		Implements the numerical solution to Eq. (12) in Ma et al. 2005.
		"""
		# 1. Calculate RHS constant C based on detector error
		h2_e = binary_entropy(e_detector)
		numerator = error_correction_f * h2_e
		denominator = 1.0 - h2_e
		
		if denominator <= 0:
		logger.warning(
		"Denominator in optimal mu calc is <= 0 "
		"(Error rate too high). Defaulting to 0.48.")
		return 0.48
		
		C = numerator / denominator
	\end{lstlisting}
\end{LTR}

در خطوط ابتدایی، ابتدا آنتروپی باینری خطای آشکارساز $H_2(e_{\text{det}})$ با تابع \lr{\texttt{binary\_entropy}} از ماژول \lr{\texttt{qkd.utils.math}} محاسبه می‌شود. سپس صورت و مخرج ثابت $C = f_{\text{EC}} H_2(e_{\text{det}}) / (1 - H_2(e_{\text{det}}))$ محاسبه می‌شوند. اگر مخرج غیرمثبت باشد (یعنی $H_2(e_{\text{det}}) \ge 1$ که فقط در $e_{\text{det}} = 0.5$ رخ می‌دهد و نشان‌دهندۀ خطای بسیار بالای آشکارساز است)، یک هشدار در لاگ ثبت می‌شود و مقدار پیش‌فرض \lr{0.48} برگردانده می‌شود. این رفتار \lr{fail-safe} تضمین می‌کند که تابع حتی در شرایط غیرفیزیکی هم یک خروجی معتبر برمی‌گرداند. سپس ثابت $C$ محاسبه می‌شود که در معادلۀ $(1 - \mu) e^{-\mu} = C$ استفاده می‌شود.

سپس تابع هدف و جواب آن به صورت زیر تعریف می‌شود:

\begin{LTR}
	\begin{lstlisting}[language=Python]
		# 2. Define function to find root: (1-mu)e^-mu - C = 0
		def objective(mu):
		return (1.0 - mu) * math.exp(-mu) - C
		
		# Define dynamic fallback based on Paper Section 3.1
		fallback_mu = 0.54 if error_correction_f <= 1.05 else 0.48
		
		# 3. Solve for mu in (0, 1]
		try:
		val_low = objective(0.001)
		val_high = objective(0.999)
		
		if val_low * val_high < 0:
		mu_opt = brentq(objective, 0.001, 0.999)
		return float(mu_opt)
		else:
		logger.warning(
		f"Optimal mu equation has no solution in [0,1]. "
		f"Falling back to {fallback_mu}.")
		return fallback_mu
		except Exception as e:
		logger.warning(
		f"Numerical solver failed for optimal mu: {e}. "
		f"Defaulting to {fallback_mu}.")
		return fallback_mu
	\end{lstlisting}
\end{LTR}

در این بخش، ابتدا تابع هدف به صورت $f(\mu) = (1 - \mu) e^{-\mu} - C$ تعریف می‌شود. این تابع در $\mu = 0$ برابر $1 - C$ (مثبت برای $C < 1$) و در $\mu \to \infty$ به $-C$ (منفی) میل می‌کند. بنابراین یک ریشه در بازۀ $(0, \infty)$ وجود دارد. برای اطمینان از وجود ریشه در بازۀ $[0.001, 0.999]$، مقدار تابع در دو انتها بررسی می‌شود: اگر $f(0.001) \cdot f(0.999) < 0$ (یعنی علامت‌های متضاد)، طبق \lr{Intermediate Value Theorem} یک ریشه در این بازه وجود دارد و \lr{\texttt{brentq}} می‌تواند آن را پیدا کند. در غیر این صورت، یک هشدار ثبت می‌شود و مقدار پیش‌فرض \lr{\texttt{fallback\_mu}} (که بر اساس $f_{\text{EC}}$ انتخاب می‌شود) برگردانده می‌شود. این \lr{fallback} بر اساس بخش 3.1 مقالۀ \lr{Ma 2005} انتخاب شده است: برای $f_{\text{EC}} \le 1.05$ مقدار \lr{0.54} و برای $f_{\text{EC}} > 1.05$ مقدار \lr{0.48} پیشنهاد شده است. در صورت بروز هرگونه استثنا در \lr{\texttt{brentq}}، بلوک \lr{\texttt{try/except}} آن را می‌گیرد و مقدار پیش‌فرض برمی‌گرداند.

\subsection{بهینه‌سازی $\nu$ با روش \lr{Brent}}

تابع \lr{\texttt{optimize\_nu\_ma2005}} بهینه‌سازی $\nu$ را برای $N$ متناهی پیاده‌سازی می‌کند:

\begin{LTR}
	\begin{lstlisting}[language=Python]
		def optimize_nu_ma2005(
		mu: float,
		total_pulses_N: float,
		cost_function_callback: Callable[[float], float],
		distance_km: float = 0.0
		) -> float:
		"""
		Optimizes the weak decoy intensity (nu) for a fixed data size N.
		Based on Ma et al. 2005 Section 4.3.
		"""
		lower_bound = 0.001
		upper_bound = mu * 0.9
		
		# Robust 1D scalar minimization (Brent's method)
		res = minimize_scalar(
		cost_function_callback,
		bounds=(lower_bound, upper_bound),
		method='bounded',
		options={'xatol': 1e-4}
		)
		
		if res.success:
		return float(res.x)
		
		# Fallback if optimization fails
		return 0.1
	\end{lstlisting}
\end{LTR}

در این تابع، ابتدا کران‌های جستجو به صورت $[\nu_{\min}, \nu_{\max}] = [0.001, 0.9 \mu]$ تعریف می‌شوند. کران پایین \lr{0.001} برای جلوگیری از مقادیر خیلی کوچک (که از نظر آماری معنایی ندارند) و کران بالا $0.9 \mu$ برای حفظ قابلیت تفکیک آماری بین سیگنال و \lr{decoy} انتخاب شده‌اند. سپس \lr{\texttt{minimize\_scalar}} با روش \lr{\texttt{'bounded'}} فراخوانی می‌شود. این روش بر اساس الگوریتم \lr{Brent's bounded method} است که برای توابع یک‌بعدی هموار بهینه است و همگرایی تضمین‌شده‌ای در بازۀ مشخص‌شده دارد. گزینه \lr{\texttt{xatol=1e-4}} تلورانس توقف را روی $10^{-4}$ تنظیم می‌کند که برای کاربردهای عملی کافی است. اگر بهینه‌سازی موفق شود، مقدار بهینۀ $\nu^*$ برگردانده می‌شود. در غیر این صورت، مقدار پیش‌فرض \lr{0.1} (یک مقدار تجربی خوب برای \lr{decoy} ضعیف) برگردانده می‌شود.

\subsection{بهینه‌سازی احتمالات توزیع پالس}

تابع \lr{\texttt{optimize\_pulse\_distribution\_ma2005}} احتمالات توزیع پالس $P_s$ و $P_d$ را بهینه‌سازی می‌کند:

\begin{LTR}
	\begin{lstlisting}[language=Python]
		def optimize_pulse_distribution_ma2005(
		mu: float,
		nu: float,
		total_pulses_N: int,
		cost_function_callback: Callable[[float, float], float]
		) -> Tuple[float, float, float]:
		"""
		Optimizes the pulse probabilities (Signal, Decoy, Vacuum) for a fixed data size N.
		"""
		def objective(x):
		p_s, p_d = x
		if p_s <= 0 or p_d <= 0 or (p_s + p_d) >= 1.0:
		return 1e9  # Penalty
		return cost_function_callback(p_s, p_d)
		
		# Initial guess
		x0 = [0.66, 0.29]
		bounds = [(0.01, 0.98), (0.01, 0.98)]
		constraints = [{'type': 'ineq', 'fun': lambda x: 1.0 - (x[0] + x[1]) - 0.01}]
		
		res = minimize(objective, x0, bounds=bounds,
		constraints=constraints, method='SLSQP')
		
		if res.success:
		p_s, p_d = res.x
		p_v = 1.0 - p_s - p_d
		return p_s, p_d, p_v
		
		logger.warning("Pulse distribution optimization failed. Using defaults.")
		return 0.8, 0.1, 0.1
	\end{lstlisting}
\end{LTR}

در این تابع، ابتدا تابع هدف داخلی تعریف می‌شود که محدودیت‌های فیزیکی را بررسی می‌کند: $P_s > 0$، $P_d > 0$ و $P_s + P_d < 1$. اگر هر کدام نقض شود، جریمۀ $10^9$ برگردانده می‌شود. در غیر این صورت، \lr{callback} فراخوانی می‌شود. سپس نقطۀ شروع $x_0 = [0.66, 0.29]$ (یک توزیع متعادل با $P_v \approx 0.05$) و کران‌های $[0.01, 0.98]$ برای هر دو پارامتر تعریف می‌شوند. یک محدودیت نامساوی نیز اضافه می‌شود: $P_s + P_d \le 0.99$ که فضایی برای $P_v \ge 0.01$ باقی می‌گذارد. سپس \lr{\texttt{minimize}} با روش \lr{\texttt{SLSQP}} فراخوانی می‌شود. این روش برای مسائل بهینه‌سازی مقید مناسب است و از محدودیت‌های نامساوی پشتیبانی می‌کند. اگر بهینه‌سازی موفق شود، $P_s^*$، $P_d^*$ و $P_v^* = 1 - P_s^* - P_d^*$ برگردانده می‌شوند. در غیر این صورت، مقادیر پیش‌فرض $(0.8, 0.1, 0.1)$ برگردانده می‌شوند که یک توزیع متعادل برای پروتکل \lr{Vacuum + Weak} است.

\subsection{بهینه‌سازی توامان سه پارامتر}

تابع \lr{\texttt{optimize\_simultaneous\_ma2005}} بهینه‌سازی توامان $P_s$، $P_d$ و $\nu$ را با استفاده از کلاس \lr{\texttt{Ma2005GlobalOptimizer}} پیاده‌سازی می‌کند:

\begin{LTR}
	\begin{lstlisting}[language=Python]
		def optimize_simultaneous_ma2005(
		mu_target: float,
		total_pulses_N: int,
		cost_function_callback: Callable[[float, float, float], float]
		) -> Tuple[float, float, float, float]:
		"""
		Simultaneously optimizes pulse probabilities (P_s, P_d) and decoy intensity (nu).
		"""
		optimizer = Ma2005GlobalOptimizer(
		mu_target, total_pulses_N, cost_function_callback)
		
		# Initial guess: Start with standard values
		x0 = [0.66, 0.29, 0.127]
		
		# Bounds
		bounds = [
		(0.1, 0.95), (0.01, 0.8), (0.001, mu_target * 0.99)
		]
		
		# Relax vacuum constraint to allow P_v -> 0
		constraints = [{'type': 'ineq',
			'fun': lambda x: 1.0 - (x[0] + x[1]) - 1e-9}]
		
		# Use SLSQP with the Jacobian method implemented in the optimizer class
		res = minimize(
		optimizer.objective,
		x0,
		jac=optimizer.gradient,
		bounds=bounds,
		constraints=constraints,
		method='SLSQP',
		tol=1e-5
		)
		
		if res.success:
		p_s_opt, p_d_opt, nu_opt = res.x
		p_v_opt = 1.0 - p_s_opt - p_d_opt
		return p_s_opt, p_d_opt, p_v_opt, nu_opt
		else:
		# If optimization fails, try fallback or return initial guess
		return 0.66, 0.29, 0.05, 0.127
	\end{lstlisting}
\end{LTR}

در این تابع، ابتدا یک نمونۀ \lr{\texttt{Ma2005GlobalOptimizer}} با \lr{\texttt{mu\_target}}، \lr{\texttt{total\_pulses\_N}} و \lr{\texttt{cost\_function\_callback}} ساخته می‌شود. سپس نقطۀ شروع $x_0 = [0.66, 0.29, 0.127]$ (مقادیر استاندارد برای پروتکل \lr{Vacuum + Weak}) تعریف می‌شود. کران‌ها به صورت $[0.1, 0.95]$ برای $P_s$، $[0.01, 0.8]$ برای $P_d$ و $[0.001, 0.99 \mu]$ برای $\nu$ تعریف می‌شوند. یک محدودیت نامساوی نیز اضافه می‌شود: $P_s + P_d \le 1 - 10^{-9}$ که اجازه می‌دهد $P_v$ به صفر میل کند (برای حالتی که \lr{decoy} خلأ لازم نیست). سپس \lr{\texttt{minimize}} با روش \lr{\texttt{SLSQP}} فراخوانی می‌شود و هم تابع هدف (\lr{\texttt{optimizer.objective}}) و هم گرادیان (\lr{\texttt{optimizer.gradient}}) به آن پاس داده می‌شوند. تلورانس توقف $10^{-5}$ تنظیم می‌شود. اگر بهینه‌سازی موفق شود، $P_s^*$، $P_d^*$، $P_v^*$ و $\nu^*$ برگردانده می‌شوند. در غیر این صورت، مقادیر پیش‌فرض $(0.66, 0.29, 0.05, 0.127)$ برگردانده می‌شوند.

\subsection{بهینه‌سازی کلید متناهی \lr{Lim 2014}}

تابع \lr{\texttt{optimize\_lim2014\_finite\_key}} بهینه‌سازی پنج پارامتر چارچوب \lr{Lim 2014} را پیاده‌سازی می‌کند:

\begin{LTR}
	\begin{lstlisting}[language=Python]
		def optimize_lim2014_finite_key(
		params: QKDParams,
		dist_km: float,
		cost_callback: Callable[[float, float, float, float, float], float]
		) -> Tuple[float, float, float, float, float]:
		"""
		Numerically optimizes the secret key rate R = l/N over the free parameters
		{q_x, p_mu1, p_mu2, mu1, mu2} as specified in Lim et al. (2014) Section IV.
		"""
		x0 = [0.9, 0.6, 0.2, 0.5, 0.1]
		
		bounds = [
		(0.5, 0.99),  # q_x
		(0.1, 0.8),   # p_s
		(0.05, 0.5),  # p_d
		(0.2, 0.8),   # mu
		(0.01, 0.2)   # nu
		]
		
		def objective(x):
		q_x, p_s, p_d, mu, nu = x
		if p_s + p_d >= 0.99: return 1e9
		if nu >= mu * 0.9: return 1e9
		if mu <= 0 or nu <= 0: return 1e9
		return cost_callback(q_x, p_s, p_d, mu, nu)
		
		# Run optimization (Nelder-Mead is robust for non-smooth finite-key landscapes)
		res = minimize(
		objective,
		x0,
		bounds=bounds,
		method='Nelder-Mead',
		options={'maxiter': 200, 'xatol': 1e-3}
		)
		
		if res.success:
		return tuple(res.x)
		
		# Fallback to initial guess if optimization fails
		logger.warning("Lim2014 parameter optimization failed. Using defaults.")
		return tuple(x0)
	\end{lstlisting}
\end{LTR}

در این تابع، نقطۀ شروع $x_0 = [0.9, 0.6, 0.2, 0.5, 0.1]$ برای پنج پارامتر $q_x$، $P_s$، $P_d$، $\mu$ و $\nu$ تعریف می‌شود. کران‌ها به صورت $[0.5, 0.99]$ برای $q_x$ (احتمال پایۀ $X$، باید بیشتر از \lr{0.5} باشد زیرا پایۀ تولیدکنندۀ کلید است)، $[0.1, 0.8]$ برای $P_s$، $[0.05, 0.5]$ برای $P_d$، $[0.2, 0.8]$ برای $\mu$ و $[0.01, 0.2]$ برای $\nu$ تعریف می‌شوند. تابع هدف داخلی محدودیت‌های فیزیکی را بررسی می‌کند: $P_s + P_d < 0.99$، $\nu < 0.9 \mu$ و $\mu, \nu > 0$. اگر هر کدام نقض شود، جریمۀ $10^9$ برگردانده می‌شود. سپس \lr{\texttt{minimize}} با روش \lr{\texttt{Nelder-Mead}} فراخوانی می‌شود. این روش بدون گرادیان است و برای توابع ناهموار مقاوم‌تر از روش‌های مبتنی بر گرادیان است. حداکثر تعداد تکرارها 200 و تلورانس $10^{-3}$ تنظیم می‌شود. اگر بهینه‌سازی موفق شود، تاپل پنج عنصری $(q_x^*, P_s^*, P_d^*, \mu^*, \nu^*)$ برگردانده می‌شود. در غیر این صورت، مقادیر پیش‌فرض (همان $x_0$) برگردانده می‌شوند.

\subsection{بهینه‌سازی زاویۀ چرخش در \lr{Chapman 2018}}

تابع \lr{\texttt{optimize\_rotation\_angle\_chapman}} بهینه‌سازی زاویۀ چرخش $\theta$ برای \lr{Coherent Scheme} در \lr{Chapman 2018} را پیاده‌سازی می‌کند:

\begin{LTR}
	\begin{lstlisting}[language=Python]
		def optimize_rotation_angle_chapman(gamma: float, M: int = 1) -> float:
		"""
		Finds the optimal encoding rotation angle theta_gamma for the Coherent Scheme
		in Chapman et al. (2018).
		
		We maximize P_succ = 0.5 * (1 + D(rho_0, rho_1)) where rho_i are outputs of ADC.
		"""
		# Kraus ops for ADC(gamma)
		E0 = np.array([[1, 0], [0, math.sqrt(1-gamma)]])
		E1 = np.array([[0, math.sqrt(gamma)], [0, 0]])
		
		# Rotation Y
		Ry = lambda theta: np.array([
		[math.cos(theta/2), -math.sin(theta/2)],
		[math.sin(theta/2),  math.cos(theta/2)]])
		
		def get_rho_out(state_in):
		return E0 @ state_in @ E0.conj().T + E1 @ state_in @ E1.conj().T
		
		def objective(theta):
		# Rotate |0> and |1> input states by theta
		U = Ry(theta)
		state0 = U @ np.array([[1,0],[0,0]]) @ U.conj().T
		state1 = U @ np.array([[0,0],[0,1]]) @ U.conj().T
		
		rho0_out = get_rho_out(state0)
		rho1_out = get_rho_out(state1)
		
		# Maximize Trace Distance for M=1
		diff = rho0_out - rho1_out
		evals = np.linalg.eigvalsh(diff)
		dist = 0.5 * np.sum(np.abs(evals))
		return -dist
		
		res = minimize_scalar(objective, bounds=(0, math.pi), method='bounded')
		return float(res.x)
	\end{lstlisting}
\end{LTR}

در این تابع، ابتدا عملگرهای \lr{Kraus} برای \lr{ADC} با پارامتر $\gamma$ ساخته می‌شوند: $E_0$ و $E_1$ به صورت ماتریس‌های \lr{NumPy} تعریف می‌شوند. سپس یک \lr{lambda} برای ماتریس چرخش \lr{Y} به صورت $R_y(\theta) = \begin{pmatrix} \cos(\theta/2) & -\sin(\theta/2) \\ \sin(\theta/2) & \cos(\theta/2) \end{pmatrix}$ تعریف می‌شود. تابع \lr{\texttt{get\_rho\_out}} حالت خروجی کانال را از رابطۀ $\rho_{\text{out}} = E_0 \rho_{\text{in}} E_0^\dagger + E_1 \rho_{\text{in}} E_1^\dagger$ محاسبه می‌کند. در تابع هدف، ابتدا ماتریس چرخش $U = R_y(\theta)$ ساخته می‌شود و حالت‌های ورودی $|0\rangle\langle 0|$ و $|1\rangle\langle 1|$ با این چرخش تبدیل می‌شوند: $\rho_i^{\text{in}} = U |i\rangle\langle i| U^\dagger$. سپس حالت‌های خروجی محاسبه می‌شوند و فاصلۀ \lr{trace} میان آن‌ها به صورت $D = \frac{1}{2} \sum_i |\lambda_i|$ محاسبه می‌شود که در آن $\lambda_i$ مقادیر ویژۀ $\rho_0^{\text{out}} - \rho_1^{\text{out}}$ هستند. تابع \lr{\texttt{eigvalsh}} از \lr{NumPy} برای محاسبۀ مقادیر ویژۀ یک ماتریس \lr{Hermitian} استفاده می‌شود (که سریع‌تر و پایدارتر از \lr{\texttt{eigvals}} عمومی است). چون بهینه‌سازی به دنبال بیشینه‌کردن $D$ است و \lr{\texttt{minimize\_scalar}} به دنبال کمینه‌کردن است، علامت منفی $-D$ برگردانده می‌شود. در نهایت، \lr{\texttt{minimize\_scalar}} با روش \lr{\texttt{'bounded'}} در بازۀ $[0, \pi]$ فراخوانی می‌شود و زاویۀ بهینۀ $\theta^*$ برگردانده می‌شود.

\subsection{بهینه‌سازی اکتشافی شدت‌ها}

تابع \lr{\texttt{optimize\_decoy\_intensities}} یک بهینه‌سازی اکتشافی برای شدت سیگنال و \lr{decoy} ارائه می‌کند:

\begin{LTR}
	\begin{lstlisting}[language=Python]
		def optimize_decoy_intensities(
		distance_km: float,
		fiber_loss_db_km: float,
		det_eff: float,
		dark_rate: float,
		error_correction_f: float
		) -> Tuple[float, float]:
		"""
		Heuristic optimization for signal and decoy intensities.
		Returns: (mu_opt, nu_opt)
		"""
		transmittance = 10**(-fiber_loss_db_km * distance_km / 10.0)
		
		def neg_key_rate(x):
		mu, nu = x[0], x[1]
		if mu <= nu or nu <= 0 or mu > 1.0:
		return 1.0
		
		# Heuristic simplified rate model
		Q_mu = 1 - math.exp(-transmittance * mu) + 2 * dark_rate
		E_mu = (0.5 * 2 * dark_rate) / Q_mu
		Y1 = transmittance * det_eff
		Q1 = mu * math.exp(-mu) * Y1
		
		rate = Q1 - error_correction_f * Q_mu * binary_entropy(E_mu)
		return -rate
		
		res = minimize(neg_key_rate, [0.5, 0.1],
		bounds=[(0.1, 0.9), (0.01, 0.3)])
		
		if res.success:
		return res.x[0], res.x[1]
		return 0.48, 0.05
	\end{lstlisting}
\end{LTR}

در این تابع، ابتدا عبور کانال به صورت $\eta = 10^{-\alpha L / 10}$ محاسبه می‌شود که در آن $\alpha$ ضریب افت فیبر (به \lr{dB/km}) و $L$ فاصله (به \lr{km}) است. این فرمول از قانون \lr{Beer-Lambert} ناشی می‌شود. سپس تابع هدف به صورت نرخ منفی کلید با یک مدل ساده‌شده تعریف می‌شود: $Q_\mu = 1 - e^{-\eta \mu} + 2 Y_0$ (احتمال آشکارسازی سیگنال به همراه \lr{dark count})، $E_\mu = Y_0 / Q_\mu$ (نرخ خطای سیگنال که عمدتاً از \lr{dark count} ناشی می‌شود)، $Y_1 = \eta \cdot \text{det\_eff}$ (بازده تک‌فوتونی) و $Q_1 = \mu e^{-\mu} Y_1$ (\lr{Gain} تک‌فوتونی). نرخ کلید به صورت $R = Q_1 - f_{\text{EC}} Q_\mu H_2(E_\mu)$ محاسبه می‌شود. این یک مدل اکتشافی است که فرض می‌کند $e_1 \approx 0$ و $q = 1$ (یعنی پایه‌گذاری بی‌اتلاف). سپس \lr{\texttt{minimize}} با نقطۀ شروع $[0.5, 0.1]$ و کران‌های $[0.1, 0.9]$ برای $\mu$ و $[0.01, 0.3]$ برای $\nu$ فراخوانی می‌شود. اگر بهینه‌سازی موفق شود، $(\mu^*, \nu^*)$ برگردانده می‌شود. در غیر این صورت، مقادیر پیش‌فرض $(0.48, 0.05)$ برگردانده می‌شوند.

\section{ارسال و دریافت با سایر ماژول‌ها}

\subsection{ورودی‌های ماژول}

ماژول \lr{\texttt{optimization.py}} برای انجام بهینه‌سازی‌های خود به چندین نوع ورودی نیاز دارد. این ورودی‌ها به دو دستۀ پارامترهای فیزیکی و \lr{callback}‌های تابعی تقسیم می‌شوند.

\begin{itemize}
	\item پارامترهای فیزیکی: این پارامترها مشخصات سیستم \lr{QKD} را توصیف می‌کنند و شامل موارد زیر هستند: \lr{\texttt{distance\_km}} (فاصله به کیلومتر)، \lr{\texttt{fiber\_loss\_db\_km}} (ضریب افت فیبر به \lr{dB/km}، معمولاً 0.2 برای فیبر سیلیکایی در طول موج 1550 \lr{nm})، \lr{\texttt{det\_eff}} (راندمان آشکارساز، معمولاً 0.1 تا 0.3 برای \lr{SPAD})، \lr{\texttt{dark\_rate}} (نرخ شمارش تاریک آشکارساز)، \lr{\texttt{error\_correction\_f}} (بازدهی تصحیح خطا، معمولاً 1.1 تا 1.2) و \lr{\texttt{e\_detector}} (خطای ذاتی آشکارساز). این پارامترها یا به صورت آرگومان‌های مستقیم به توابع پاس داده می‌شوند یا از شیء \lr{\texttt{QKDParams}} استخراج می‌شوند.
	\item فراخوان‌های تابعی (\lr{callback}ها): این \lr{callback}‌ها تابع هدف را برای بهینه‌سازی تعریف می‌کنند و توسط کاربر ارائه می‌شوند. امضای این \lr{callback}‌ها بسته به تابع بهینه‌سازی متفاوت است: \lr{\texttt{Callable[[float], float]}} برای \lr{\texttt{optimize\_nu\_ma2005}}، \lr{\texttt{Callable[[float, float], float]}} برای \lr{\texttt{optimize\_pulse\_distribution\_ma2005}}، \lr{\texttt{Callable[[float, float, float], float]}} برای \lr{\texttt{optimize\_simultaneous\_ma2005}} و \lr{\texttt{Callable[[float, float, float, float, float], float]}} برای \lr{\texttt{optimize\_lim2014\_finite\_key}}. این \lr{callback}‌ها معمولاً نرخ منفی کلید را محاسبه می‌کنند (زیرا حل‌گرها (\lr{solver}ها) به دنبال کمینه‌کردن هستند) اما می‌توانند هر کمیت هدف دیگری را نیز محاسبه کنند.
	\item \lr{\texttt{QKDParams}}: این شیء از ماژول \lr{\texttt{qkd.params}} وارد می‌شود و در تابع \lr{\texttt{optimize\_lim2014\_finite\_key}} استفاده می‌شود. این شیء شامل تمام پارامترهای پیکربندی سیستم است و به \lr{callback} اجازه می‌دهد به پارامترهای اثبات امنیتی (مانند $\varepsilon_{\text{sec}}$، $\varepsilon_{\text{cor}}$ و $\varepsilon_{\text{pe}}$) دسترسی داشته باشد.
	\item \lr{\texttt{gamma}} (پارامتر کانال): این پارامتر در توابع \lr{\texttt{optimize\_rotation\_angle}} و \lr{\texttt{optimize\_rotation\_angle\_chapman}} استفاده می‌شود و نشان‌دهندۀ پارامتر کانال \lr{ADC} (در \lr{Chapman 2018}) یا یک کانال عمومی است.
\end{itemize}

\subsection{خروجی‌های ماژول}

خروجی‌های ماژول بسته به تابع بهینه‌سازی متفاوت هستند:

\begin{enumerate}
	\item خروجی‌های تک‌مقداری: توابع \lr{\texttt{calculate\_optimal\_mu\_ma2005}}، \lr{\texttt{optimize\_nu\_ma2005}}، \lr{\texttt{optimize\_rotation\_angle}} و \lr{\texttt{optimize\_rotation\_angle\_chapman}} یک مقدار اعشاری واحد را برمی‌گردانند (به ترتیب $\mu^*$، $\nu^*$، $\theta^*$ عمومی و $\theta^*$ برای \lr{Chapman}).
	\item خروجی‌های دومقداری: تابع \lr{\texttt{optimize\_decoy\_intensities}} یک تاپل دو عنصری شامل $(\mu^*, \nu^*)$ را برمی‌گرداند.
	\item خروجی‌های سه‌مقداری: تابع \lr{\texttt{optimize\_pulse\_distribution\_ma2005}} (و پوشش‌دهندۀ آن \lr{\texttt{optimize\_pulse\_probabilities}}) یک تاپل سه عنصری شامل $(P_s^*, P_d^*, P_v^*)$ را برمی‌گرداند.
	\item خروجی‌های چهارمقداری: تابع \lr{\texttt{optimize\_simultaneous\_ma2005}} (و پوشش‌دهندۀ آن \lr{\texttt{optimize\_global\_ma2005}}) یک تاپل چهار عنصری شامل $(P_s^*, P_d^*, P_v^*, \nu^*)$ را برمی‌گرداند.
	\item خروجی‌های پنج‌مقداری: تابع \lr{\texttt{optimize\_lim2014\_finite\_key}} یک تاپل پنج عنصری شامل $(q_x^*, P_s^*, P_d^*, \mu^*, \nu^*)$ را برمی‌گرداند.
\end{enumerate}

در تمام موارد، اگر بهینه‌سازی شکست بخورد، مقادیر پیش‌فرض محافظه‌کارانه برگردانده می‌شوند. این مقادیر پیش‌فرض بر اساس تجربۀ عملی و مقالات مرجع انتخاب شده‌اند و برای پروتکل‌های \lr{Decoy-State BB84} معمولاً خوب عمل می‌کنند.

\subsection{وابستگی‌های داخلی و خارجی}

ماژول وابستگی‌های داخلی و خارجی متعددی دارد. وابستگی‌های داخلی شامل موارد زیر است:

\begin{itemize}
	\item \lr{\texttt{qkd.params.QKDParams}}: کلاس پارامترهای \lr{QKD} که در تابع \lr{\texttt{optimize\_lim2014\_finite\_key}} استفاده می‌شود.
	\item \lr{\texttt{qkd.utils.math.binary\_entropy}}: تابع محاسبۀ آنتروپی باینری $H_2(p) = -p \log_2 p - (1-p) \log_2 (1-p)$ که در توابع \lr{\texttt{calculate\_optimal\_mu\_ma2005}} و \lr{\texttt{optimize\_decoy\_intensities}} استفاده می‌شود.
	\item \lr{\texttt{qkd.constants.is\_valid\_probability}}: تابع اعتبارسنجی احتمال که بررسی می‌کند یک مقدار در بازۀ $[0, 1]$ قرار دارد.
	\item \lr{\texttt{qkd.datatypes.TallyCounts}}: کلاس دادهٔ (\lr{dataclass}) شمارش‌های آماری که در برخی \lr{callback}‌ها برای محاسبۀ نرخ کلید استفاده می‌شود.
\end{itemize}

وابستگی‌های خارجی شامل موارد زیر است:

\begin{itemize}
	\item \lr{\texttt{scipy.optimize.minimize}}: تابع بهینه‌سازی چندبعدی \lr{SciPy} که در توابع \lr{\texttt{optimize\_decoy\_intensities}}، \lr{\texttt{optimize\_pulse\_distribution\_ma2005}}، \lr{\texttt{optimize\_simultaneous\_ma2005}} و \lr{\texttt{optimize\_lim2014\_finite\_key}} استفاده می‌شود. این تابع از روش‌های مختلفی (مانند \lr{SLSQP} و \lr{Nelder-Mead}) پشتیبانی می‌کند.
	\item \lr{\texttt{scipy.optimize.minimize\_scalar}}: تابع بهینه‌سازی یک‌بعدی \lr{SciPy} که در توابع \lr{\texttt{optimize\_rotation\_angle}}، \lr{\texttt{optimize\_nu\_ma2005}} و \lr{\texttt{optimize\_rotation\_angle\_chapman}} استفاده می‌شود. این تابع از روش \lr{Brent's bounded} پشتیبانی می‌کند.
	\item \lr{\texttt{scipy.optimize.brentq}}: تابع یافتن ریشه \lr{SciPy} که در تابع \lr{\texttt{calculate\_optimal\_mu\_ma2005}} برای حل معادلۀ $(1-\mu)e^{-\mu} = C$ استفاده می‌شود.
	\item \lr{\texttt{numpy}}: کتابخانۀ محاسبات برداری که در تابع \lr{\texttt{optimize\_rotation\_angle\_chapman}} برای محاسبۀ ماتریس‌ها و مقادیر ویژه استفاده می‌شود.
	\item \lr{\texttt{math}}: ماژول ریاضی استاندارد \lr{Python} برای توابع پایه مانند \lr{\texttt{exp}}، \lr{\texttt{log}}، \lr{\texttt{cos}}، \lr{\texttt{sin}} و \lr{\texttt{sqrt}}.
	\item \lr{\texttt{logging}}: ماژول لاگ‌گذاری استاندارد \lr{Python} برای ثبت هشدارها و اطلاعات.
	\item \lr{\texttt{typing}}: ماژول \lr{Type Hint} برای تعریف امضای توابع.
\end{itemize}

\subsection{جریان داده در سراسر فریم‌ورک}

جریان داده کلی به این شکل است:
\begin{enumerate}
	\item لایۀ بالایی (مانند اسکریپت پیکربندی یا شبیه‌ساز) پارامترهای فیزیکی سیستم (مانند فاصله، افت فیبر، راندمان آشکارساز) را مشخص می‌کند.
	\item یک \lr{callback} تابعی تعریف می‌شود که نرخ کلید (یا هر کمیت هدف دیگر) را برای مجموعه‌ای از پارامترهای داده‌شده محاسبه می‌کند. این \lr{callback} معمولاً از ماژول‌های اثبات امنیتی (مانند \lr{\texttt{lim2014.py}}، \lr{\texttt{tight.py}} یا \lr{\texttt{ma2005.py}}) استفاده می‌کند.
	\item تابع بهینه‌سازی مناسب (مانند \lr{\texttt{optimize\_lim2014\_finite\_key}}) با پارامترهای فیزیکی و \lr{callback} فراخوانی می‌شود.
	\item تابع بهینه‌سازی، \lr{callback} را در نقاط مختلف فضای پارامتر فراخوانی می‌کند و با استفاده از روش‌های عددی (مانند \lr{SLSQP} یا \lr{Nelder-Mead}) به سمت نقطۀ بهینه حرکت می‌کند.
	\item در هر فراخوانی \lr{callback}، تابع هدف ابتدا محدودیت‌های فیزیکی را بررسی می‌کند و در صورت نقض، جریمۀ بزرگی برگردانده می‌شود.
	\item پس از همگرایی (یا رسیدن به حداکثر تعداد تکرار)، تابع بهینه‌سازی پارامترهای بهینه را برمی‌گرداند.
	\item لایۀ بالایی از این پارامترهای بهینه برای پیکربندی شبیه‌سازی نهایی استفاده می‌کند.
\end{enumerate}

این جداسازی مراحل به گونه‌ای طراحی شده که هر مرحله به صورت مستقل قابل آزمایش و بازبینی باشد. به عنوان مثال، می‌توان از \lr{callback}‌های مصنوعی با رفتار مشخص برای اعتبارسنجی توابع بهینه‌سازی استفاده کرد. این ماژول‌بندی از اصل تک‌مسئولیتی (\lr{Single Responsibility}) پیروی می‌کند و قابلیت نگهداری و توسعه کد را به طور قابل توجهی افزایش می‌دهد. علاوه بر این، الگوی \lr{callback} به کاربر اجازه می‌دهد تابع هدف خود را بدون تغییر کد بهینه‌سازی جایگزین کند. این انعطاف‌پذیری برای تحقیقات \lr{QKD} که در آن‌ها معیارهای مختلف بهینه‌سازی (نرخ کلید، \lr{QBER}، طول کلید متناهی و غیره) مورد بررسی قرار می‌گیرند، حیاتی است. طراحی مطمئن از خرابی (\lr{fail-safe}) این ماژول، آن را به انتخابی ایده‌آل برای محیط‌های عملیاتی تبدیل می‌کند که در آن‌ها پایداری و قابلیت اطمینان سیستم اهمیت زیادی دارد.
